%----------------------------------------------------------------------------
\chapter*{Összefoglaló}\addcontentsline{toc}{chapter}{Összefoglaló}
%----------------------------------------------------------------------------

A DineEase projekt egy modern asztalfoglalórendszer megvalósításával foglalkozik, amely megoldást nyújt úgy a foglalni vágyó embereknek, mint az éttermeknek a foglalások egyszerűbb kezelésére. Lehetőséget nyújt az alkalmazás a felhasználói vélemények megnyilvánulásának is, amely alapján könnyebben dönthet az éttermek között. Ugyanakkor az éttermeknek is egy olyan felületet biztosít, amelyen egyszerűen kezelheti foglalásait, vagy akár meetingjeit és eseményeit is könnyedén hirdetheti, illetve elemzéseket is végezhet a foglalások alapján. Hozzájárul az olyan éttermeknek az eléréséhez is, amelyek nagyobb eseményekre specializálódottak, így könnyebbé téve az éttermekkel való kommunikációt.

A dolgozatom során ismertettem a hasonló alkalmazásokat és összemértem azokat a DineEase alkalmazásal funkcionalitás és technológia szempontjából is. A biztonsági lépéseket is leírtam dolgozatom során, amellyel hozzájárulhattam a felhasználók adatainak védelméhez. Majd a funkcionalitásokat is tárgyaltam, amelyek megvalósításra kerültek, illetve a funkcionális és nem funkcionális rendszerkövetelményekben is megfogalmaztam, hogy mik az elvárások a rendszerrel kapcsolatban. Ezekután tárgyaltam, hogyan épül fel a rendszer, amely egy .NET keretrendszerrel készített API-ból, egy MSSQL adatbázisból és Flutter keretrendszerrel készült frontend-ből tevődik össze. Itt bemutatásra került a verziókövető rendszer és a projektmenedzsment eszköz is. A megvalósítás részben tárgyalva van, hogyan épülnek fel a különböző komponensek. Mindezek után pedig a felhasználói felület is bemutatásra került.

Összességében úgy gondolom, sikerült egy jól használható alkalmazást fejleszteni, amely könnyebbséget jelenthet a felhasználók életében. Ugyanakkor több különböző tervvel járulnék hozzá az alkalmazás további fejlesztéséhez:

\begin{itemize}
    \item \textbf{Google térkép}: ezzel hozzájárulva a könnyebb helymeghatározáshoz, esetleg legközelebbi hely szerinti kereséshez.

    \item \textbf{Kép feltöltés}: a menü, esemény és a felhasználó táblában helyet biztosítani a képeknek, amelyek által jobb élményt nyújthatunk a felhasználóknak.

    \item \textbf{Meghívók küldése}: a foglalásban megemlítve azokat a felhasználókat, akikkel tervez az étterembe menni, meghívók küldődnének ki, hogy emlékeztessék őket.

    \item \textbf{Értesítések}: a foglalás elfogadásáról vagy elutasításáról, ajánlott eseményekről, esetleg a foglalás közeledtéről szóló értesítések küldése.
    
\end{itemize}


A részletes munkafolyamat elérhető a következő GitHub és Trello linken: 
\begin{itemize}
    \item \href{https://github.com/aksza/DineEase}{https://github.com/aksza/DineEase}
    \item \href{https://trello.com/b/FlzsShfn/dineease}{https://trello.com/b/FlzsShfn/dineease}
\end{itemize}